\assignq{3}{3}{Hydrogenic static polarisability from perturbation theory}{2}
\subsection*{(a) Approximate the \(1s\) polarizability}
For a uniform electric field in the \(z\) direction,
\(H'=-\mathbf d\cdot\boldsymbol{\mathcal E}
=e z\mathcal E\). The \(1s\) state has even
parity, so \(\langle1s|z|1s\rangle=0\):
its first-order Stark shift vanishes.
For a nondegenerate level, second-order
perturbation theory gives
\[
\Delta E^{(2)}=\mathcal E^2
\sum_{m\ne1}\frac{|\langle m|ez|1\rangle|^2}{E_1-E_m}
=-\tfrac12\alpha\mathcal E^2.
\]
In atomic units, therefore,
\[
\alpha=2\sum_{m\ne1}
\frac{|\langle m|z|1s\rangle|^2}{E_m-E_1}.
\]
Approximate every denominator by
\(E_2-E_1=3Z^2/8\) Hartree.
Closure, \(\sum_m\ket m\bra m=I\), gives
\[
\sum_{m\ne1}|\langle m|z|1s\rangle|^2
=\langle1s|z^2|1s\rangle
-|\langle1s|z|1s\rangle|^2
=\langle z^2\rangle.
\]
Spherical symmetry implies
\(\langle z^2\rangle=\langle r^2\rangle/3\).
For hydrogenic \(1s\),
\(\langle r^2\rangle=3/Z^2\), hence
\[
\boxed{\alpha_{\rm approx}=
\frac{2(1/Z^2)}{3Z^2/8}
=\frac{16}{3Z^4}\ a_0^3.}
\]

\subsection*{(b) Hydrogen: numerical comparison}
For \(Z=1\), \(\alpha_{\rm approx}=16/3=5.3333\ a_0^3\).
With \(a_0=5.2917721\times10^{-9}\,\mathrm{cm}\),
\[
\boxed{\alpha_{\rm approx}(H)
=5.3333\,a_0^3
=7.903\times10^{-25}\,\mathrm{cm}^3}.
\]
This is the standard \emph{polarizability-volume}
convention; the SI dipole polarizability is
\(4\pi\epsilon_0\) times this volume.
The exact hydrogen value is \(9a_0^3/2=4.5a_0^3\).
The approximate denominator overestimates it by
\[
\boxed{\frac{5.3333-4.5}{4.5}\times100\%
=18.52\%.}
\]

\subsection*{(c) Hydrogen-like rubidium}
For a nucleus of \(Z=37\) with \emph{one}
hydrogenic \(1s\) electron:
\[
\boxed{\alpha(\mathrm{Rb}^{36+})\approx
\frac{16}{3(37)^4}\,a_0^3
=2.846\times10^{-6}\,a_0^3
=4.217\times10^{-31}\,\mathrm{cm}^3}.
\]
Neutral Rb has outer \(5s\) valence electrons
which are far more weakly bound and more
spatially extended; its ground-state
polarisability (\(\sim319\,a_0^3\) as supplied)
is about \(1.12\times10^8\) times larger.
\methodnote{Parity makes the first-order
shift vanish. Closure converts a difficult
sum to a ground-state expectation value;
the approximate energy denominator is why
the result is not exact.}
